\section{（诉讼）调解}
\begin{enumerate}
  \item 法院组织当事人就纠纷的解决达成协议 \yjt 调解协议（把当事人说的写下来） \yjt 调解书。
  \item 以不公开（过程和内容均）为原则。当事人处分自已私权利的行为，不需要接受舆论监督。
  \item 有一定的程序上的要求。
  \item 适用范围。
        \begin{enumerate}
          \item 简易程序审理。涉及到家庭邻里、争议小的案件。离婚诉讼应当先行调解。
          \item 诉讼程序可以调，其他程序不能调。\al{支付令、执行、公示催告、特别程序 \cuo}。
          \item 身份关系确认案件不能调。
        \end{enumerate}
  \item 调解协议。
        \begin{enumerate}
          \item 尊重意思自治。可超请求、可约责任、自愿合法、不可再诉。\al{甲乙离婚，主张 ABC 请求。调解 \yjt 多了个 D 请求、债务甲承担 \dui}。
          \item 约定一方不履行协议，另一方可以请求法院作出裁判的条款，不予准许。
                \begin{enumerate}
                  \item 本来时你俩的事，去起诉就变成了法院的事。
                  \item 一事不再理。通过调解方式，法院已经对纠纷进行了实体解决。
                \end{enumerate}
        \end{enumerate}
  \item 调解书。
        \begin{enumerate}
          \item 原则应当根据调解协议制作调解书。例外。直接以调解协议结案，就这 4 种法定情形。
                \begin{enumerate}
                  \item 当庭兑现、调解和好的离婚、维持收养关系 \yjt 日后又出了矛盾，对着调解书越想越气。
                  \item 双方当事人、审判员、书记员在调解协议、笔录上签字盖章生效。这个只有法院调解有！！
                \end{enumerate}
          \item 效力。
                \begin{enumerate}
                  \item 双方签收（含送达）后生效。签收后不准反悔。（有给付内容的）可以申请强制执行。\al{都有执行力 \cuo}。
                  \item 需要无独三担责的，调解书签收后生效。
                \end{enumerate}
          \item 担保人不签收不影响调解书的效力/生效。担保条件成就时，调解书对担保人产生约束。
          \item 错误处理。
                \begin{enumerate}
                  \item 调解书和调解协议内容不一致，重新补正即可。\al{申请再审 \cuo}。
                  \item 确有错误 \yjt 再审。\al{撤销、确认无效 \cuo}。
                        不要拿民法上的合同类推。撤销一份生效的法律文书只有三撤和再审！！
                        \al{甲称调解协书是受\xhx{胁迫}所签。可以\xhx{申请撤销}调解书、请求\xhx{确认}调解书\xhx{无效}}。
                \end{enumerate}
        \end{enumerate}
  \item 和解。自行互谅达成协议。实践很难区分调解、和解，且意义不大。
        \begin{enumerate}
          \item 结案方式。\yi 可以申请法院根据和解协议制作调解书。\al{制作判决书 \cuo}。\er 撤诉。
          \item 和解协议不具有强制执行力，不履行可就和解协议再起诉。
        \end{enumerate}
  \item 不允许据调解、和解协议制作判决书。2 个例外。应当制作调解书。
        \begin{enumerate}
          \item 无民事行为能力人的离婚。
          \item 涉外民事案件。有可能要拿去申请承认和执行。调解是中国特色。
        \end{enumerate}
  \item 和解、调解履行、救济问题。
        \begin{enumerate}
          \item 法院调解。一肯一否。能执行，否起诉
                \begin{enumerate}
                  \item 不可以就调解协议起诉 \yjt 纠纷已实体处理，重复起诉。
                  \item 有强制执行力，可以申请执行。
                  \item 二审、再审的调解多一个一审判决视为撤销的功能。
                \end{enumerate}
          \item 当事人和解、调解。简化记。双重否定。
                \begin{enumerate}
                  \item 不具有强制执行力，不能申请执行。
                  \item 不履行只能就和解/调解协议起诉。
                \end{enumerate}
          \item 确认调解协议效力。\yi 提异议。\er 重新和解、调解、仲裁、起诉。
          \item 执行和解。一肯一否。能起诉，否执行。它能起诉是例外。
          \item 仲裁调解。
        \end{enumerate}
\end{enumerate}
